\documentclass[12pt]{article}

\usepackage[margin=1in]{geometry}
\usepackage{hyperref}
\usepackage{booktabs}
\usepackage{listings}
\usepackage{xcolor}
\usepackage{enumitem}
\usepackage{parskip}

% Code listing style
\lstset{
    basicstyle=\ttfamily\small,
    backgroundcolor=\color{gray!10},
    frame=single,
    breaklines=true,
    commentstyle=\color{gray},
}

\title{PSID Family Structures}
\author{Sarah Sullivan \\ University of Michigan}
\date{Created: March 26, 2026 \\ Last Updated: March 30, 2026}

\begin{document}

\maketitle

\noindent\rule{\textwidth}{0.4pt}

% -------------------------------------------------------------------
\section*{Overview}
% -------------------------------------------------------------------

This repository contains code for the 2026 family structure project using the \textbf{Panel Study of Income Dynamics (PSID)}, a longitudinal household survey collected annually (1968--1997) and biennially (1999--2023) by the University of Michigan.

\noindent\rule{\textwidth}{0.4pt}

% -------------------------------------------------------------------
\section*{Software}
% -------------------------------------------------------------------

\begin{center}
\begin{tabular}{ll}
\toprule
\textbf{Tool} & \textbf{Purpose} \\
\midrule
\textbf{Stata}        & Data cleaning, construction, and regression analysis (\texttt{.do} files) \\
\textbf{Python}       & Data wrangling and exploratory analysis (\texttt{.ipynb} notebooks) \\
\textbf{\LaTeX}       & Tables and documents (\texttt{.tex} / \texttt{.pdf}) \\
\textbf{Git / GitHub} & Version control \\
\bottomrule
\end{tabular}
\end{center}

\begin{quote}
Auxiliary \LaTeX{} files (\texttt{.aux}, \texttt{.log}, \texttt{.toc}, etc.) are automatically cleaned after each compilation via the VS Code LaTeX Workshop setting \texttt{latex-workshop.latex.autoClean.run: ``onBuilt''}. They are also excluded from the repository via \texttt{.gitignore}.
\end{quote}

\noindent\rule{\textwidth}{0.4pt}

% -------------------------------------------------------------------
\section*{Repository Structure}
% -------------------------------------------------------------------

\begin{center}
\begin{tabular}{ll}
\toprule
\textbf{File / Folder} & \textbf{Description} \\
\midrule
\texttt{00\_globals.do}         & Global macros and paths \\
\texttt{00\_fims\_*.do} / \texttt{.dta} & Family matching \\
\texttt{00\_RENAMING\_v1.do}    & Variable renaming script \\
\texttt{\_psid}                 & Main data wrangling file \\
\texttt{A1\_vs/}                & Archive Folder \\
\texttt{git-workflow-guide.md}  & Git command reference (Markdown) \\
\texttt{git-workflow-guide.tex} & Git command reference (\LaTeX/PDF) \\
\texttt{.gitignore}             & Excluded file types \\
\texttt{\_archive/}             & Older script versions \\
\bottomrule
\end{tabular}
\end{center}

Raw PSID family-year data folders (\texttt{fam1968/} through \texttt{fam2023er/}) and individual-level data (\texttt{ind2023er/}) are stored locally but excluded from the repository via \texttt{.gitignore}. For details on downloading raw PSID data, see \href{https://psidonline.isr.umich.edu/}{psid.org}.

\noindent\rule{\textwidth}{0.4pt}

% -------------------------------------------------------------------
\section*{Key Scripts}
% -------------------------------------------------------------------

\begin{center}
\begin{tabular}{ll}
\toprule
\textbf{File} & \textbf{Description} \\
\midrule
\texttt{00\_globals.do}     & Sets global paths and macros used across all scripts \\
\texttt{00\_fims\_pars.do}  & Cleans PSID parent file \\
\texttt{00\_fims\_gpars.do} & Cleans PSID grandparent file \\
\texttt{00\_fims\_sib.do}   & Cleans PSID sibling file \\
\texttt{\_psid}             & Constructs the complete analysis panel \\
\bottomrule
\end{tabular}
\end{center}

\noindent\rule{\textwidth}{0.4pt}

% -------------------------------------------------------------------
\section*{Git Workflow}
% -------------------------------------------------------------------

See \href{git-workflow-guide.md}{\texttt{git-workflow-guide.md}} (or the compiled \href{git-workflow-guide.pdf}{\texttt{git-workflow-guide.pdf}}, also available in \LaTeX{}) for a full reference of Git commands used in this project, including:

\begin{itemize}
    \item Everyday push workflow
    \item First-time repository setup
    \item \texttt{.gitignore} management
    \item Troubleshooting rejected pushes
\end{itemize}

\textbf{Quick reference --- everyday push:}

\begin{lstlisting}[language=bash]
git pull                         # get latest changes from GitHub
git add .                        # stage all changes
git commit -m "describe changes"
git push                         # upload to GitHub
\end{lstlisting}

\noindent\rule{\textwidth}{0.4pt}

% -------------------------------------------------------------------
\section*{Data Source}
% -------------------------------------------------------------------

\noindent{\setlength{\hangindent}{2em}
Panel Study of Income Dynamics, public use dataset. Produced and distributed by the Survey Research Center, Institute for Social Research, University of Michigan, Ann Arbor, MI (2025).
\par}

\bigskip\noindent\rule{\textwidth}{0.4pt}\bigskip

% -------------------------------------------------------------------
\section*{\texttt{Variables}}
% -------------------------------------------------------------------
* See excel sheet "Variables.xlsx"

\newpage

% -------------------------------------------------------------------
\section*{\texttt{00\_RENAMING\_v1.do} Documentation}
% -------------------------------------------------------------------

The file \texttt{00\_RENAMING\_v1.do} inputs each PSID year file, retains only variables of interest, and standardizes these variable names across different waves.

Comparable variables between years are identified using \texttt{psid.xlsx} and yearly codebooks, all available for download at \href{https://psidonline.isr.umich.edu/}{psid.org}.

\subsection*{01. Identity}
\begin{enumerate}[label=\alph*.]
    \item Subtask a
    \item Subtask b
    \item Subtask c
    \item Subtask d
    \item Subtask e
\end{enumerate}

\subsection*{02. Task Name}
\begin{enumerate}[label=\alph*.]
    \item Subtask a
    \item Subtask b
    \item Subtask c
    \item Subtask d
    \item Subtask e
\end{enumerate}

\subsection*{03. Task Name}
\begin{enumerate}[label=\alph*.]
    \item Subtask a
    \item Subtask b
    \item Subtask c
    \item Subtask d
    \item Subtask e
\end{enumerate}

\subsection*{04. Task Name}
\begin{enumerate}[label=\alph*.]
    \item Subtask a
    \item Subtask b
    \item Subtask c
    \item Subtask d
    \item Subtask e
\end{enumerate}

\subsection*{05. Task Name}
\begin{enumerate}[label=\alph*.]
    \item Subtask a
    \item Subtask b
    \item Subtask c
    \item Subtask d
    \item Subtask e
\end{enumerate}

\subsection*{06. Task Name}
\begin{enumerate}[label=\alph*.]
    \item Subtask a
    \item Subtask b
    \item Subtask c
    \item Subtask d
    \item Subtask e
\end{enumerate}

\subsection*{07. Task Name}
\begin{enumerate}[label=\alph*.]
    \item Subtask a
    \item Subtask b
    \item Subtask c
    \item Subtask d
    \item Subtask e
\end{enumerate}

\subsection*{08. Task Name}
\begin{enumerate}[label=\alph*.]
    \item Subtask a
    \item Subtask b
    \item Subtask c
    \item Subtask d
    \item Subtask e
\end{enumerate}

\subsection*{09. Task Name}
\begin{enumerate}[label=\alph*.]
    \item Subtask a
    \item Subtask b
    \item Subtask c
    \item Subtask d
    \item Subtask e
\end{enumerate}

\subsection*{10. Task Name}
\begin{enumerate}[label=\alph*.]
    \item Subtask a
    \item Subtask b
    \item Subtask c
    \item Subtask d
    \item Subtask e
\end{enumerate}

\subsection*{11. Task Name}
\begin{enumerate}[label=\alph*.]
    \item Subtask a
    \item Subtask b
    \item Subtask c
    \item Subtask d
    \item Subtask e
\end{enumerate}

\subsection*{12. Task Name}
\begin{enumerate}[label=\alph*.]
    \item Subtask a
    \item Subtask b
    \item Subtask c
    \item Subtask d
    \item Subtask e
\end{enumerate}

\subsection*{13. Task Name}
\begin{enumerate}[label=\alph*.]
    \item Subtask a
    \item Subtask b
    \item Subtask c
    \item Subtask d
    \item Subtask e
\end{enumerate}

\subsection*{14. Task Name}
\begin{enumerate}[label=\alph*.]
    \item Subtask a
    \item Subtask b
    \item Subtask c
    \item Subtask d
    \item Subtask e
\end{enumerate}

\subsection*{15. Task Name}
\begin{enumerate}[label=\alph*.]
    \item Subtask a
    \item Subtask b
    \item Subtask c
    \item Subtask d
    \item Subtask e
\end{enumerate}

\subsection*{16. Task Name}
\begin{enumerate}[label=\alph*.]
    \item Subtask a
    \item Subtask b
    \item Subtask c
    \item Subtask d
    \item Subtask e
\end{enumerate}

\subsection*{17. Task Name}
\begin{enumerate}[label=\alph*.]
    \item Subtask a
    \item Subtask b
    \item Subtask c
    \item Subtask d
    \item Subtask e
\end{enumerate}

\subsection*{18. Task Name}
\begin{enumerate}[label=\alph*.]
    \item Subtask a
    \item Subtask b
    \item Subtask c
    \item Subtask d
    \item Subtask e
\end{enumerate}

\subsection*{19. Task Name}
\begin{enumerate}[label=\alph*.]
    \item Subtask a
    \item Subtask b
    \item Subtask c
    \item Subtask d
    \item Subtask e
\end{enumerate}

\subsection*{20. Task Name}
\begin{enumerate}[label=\alph*.]
    \item Subtask a
    \item Subtask b
    \item Subtask c
    \item Subtask d
    \item Subtask e
\end{enumerate}







\bigskip\noindent\rule{\textwidth}{0.4pt}\bigskip

% -------------------------------------------------------------------
\section*{\texttt{\_psid} Documentation}
% -------------------------------------------------------------------

\subsection*{00.\ Program Set Up}

\subsection*{Part I}

\begin{enumerate}
    \item For survey years, 1968--2023
    \begin{enumerate}[label=\arabic*.]
        \item Import year data file (\texttt{``famYYYY.dta''} or \texttt{``famYYYYer.dta''})
        \item Rename family ID variable \texttt{``fam''}
        \item Generate variable for year
        \item Pass through \texttt{``00\_RENAMING\_V1.do''}
    \end{enumerate}
    \item Merge together all cleaned surveys for survey years 1968--2023.
\end{enumerate}

\subsection*{Part 2}

\subsection*{Part 3}

\end{document}
